% =========================================================
% 5. Expected Results
% =========================================================
\section*{五、預期結果}
\addcontentsline{toc}{section}{五、預期結果}

\subsection*{(一) 預期可量化成果}
\begin{enumerate}
  \item \textbf{跨會議邏輯一致性提升}：相較於 Baseline，具 Dependency Graph 的系統預期在邏輯矛盾率下降、依賴鏈引用完整性提升、TODO 保留率提升等指標上達到顯著改善，並能以圖中的衝突邊作為依據，輸出可檢核的錯誤定位結果 \citep{a2602_05665v1_graphagentmemory,a2505_16067v2_memorymanagementimpacts}。
  \item \textbf{Persona Drift 顯著降低}：加入 Persona Vector Monitoring + Steering 後，預期一致性分數 $S$ 的跌破閾值頻率降低、偏移後恢復速度提升，並可在壓力情境（受試者誤導/迎合誘導）下保持教授式指導風格的穩定 \citep{a2507_21509v3_personavectors,turner-etal-2023-activation-addition}。
  \item \textbf{真人指導品質提升}：在 MCA-21 與教學能力評測面向上，同時納入 Memory 與 Persona 時，預期可提升評分表現，特別是在「回饋可行性、方法論嚴謹、錯誤診斷、跨回合一致性」等面向 \citep{hyun2022mca21,maurya-etal-2025-unifying}。
\end{enumerate}

\subsection*{(二) 學術貢獻}
\begin{itemize}
  \item \textbf{提出「長期指導一致性」的系統化問題定義與可操作架構}：將跨會議一致性拆分為「研究脈絡一致性（圖式記憶可檢核）」與「人設一致性（表徵可監控可干預）」兩條主線，以補足現有教育導師研究偏重短程評測的不足 \citep{tack-etal-2023-bea,maurya-etal-2025-unifying}。
  \item \textbf{整合圖式長期記憶與推理時人設控制的系統化方法}：以多層次記憶協調降低長程邏輯矛盾，同時以 persona 向量在推理時提供可量測、可校正的穩定性機制 \citep{a2310_08560v2_memgpt,a2602_05665v1_graphagentmemory,a2507_21509v3_personavectors}。
  \item \textbf{建立可重現的評測流程}：包含邏輯矛盾挑戰、漂移量化與真人盲測，使「一致性」由描述性主張轉為可比較的量化結果 \citep{a2402_17753v1_locomomemoryeval,hyun2022mca21}。
\end{itemize}

\subsection*{(三) 實務價值}
\begin{itemize}
  \item \textbf{低成本可部署}：採推理時干預而非 retraining，降低計算與資料需求，並提升在校園或研究室環境中落地的可行性 \citep{turner-etal-2023-activation-addition,wehner-etal-2025-representation-engineering}。
  \item \textbf{提升學術指導的可靠性與可追溯性}：以 Dependency Graph 支援決策理由追溯與矛盾定位，降低因長程錯誤累積造成的研究方向偏移風險 \citep{a2602_05665v1_graphagentmemory,liu2025agendaaware_meetingsum}。
\end{itemize}